\section{Research methodology}

\msub{1.1.1 Research design}
This study uses a quantitative, descriptive and correlational design with
secondary data \citep{pant2016}. The descriptive analysis summarizes the
levels and fiscal-year patterns of EPS, DPS, BVPS, ROA, and
$\mathit{MPS}^{post}$. Pearson correlations describe pairwise associations,
and pooled ordinary least squares (OLS) estimates the joint conditional
associations specified in Eq.~\ref{eq:main}. The intended sample is five bank
series observed in each of six fiscal years, FY 2076/77--2081/82, for a target
of 30 bank-year observations, subject to verification of the underlying
reports and trading records. Completeness is not assumed before data
collection. Any unavailable or unverifiable observation will be identified
and reported, and will not be imputed silently. If one or more observations
cannot be established, the realized sample and implications for estimation
will be reported. The design specifies six annual observations for each bank
series, but four series cross a merger boundary: pre-merger observations refer
to the acquiring institution and later observations refer to the consolidated
institution. These series are not unchanged legal entities throughout the
period. The OLS analysis is exploratory; its conventional significance tests
do not account for possible within-series dependence and will not be treated
as confirmatory evidence or as causal effects.

\msub{1.1.2 Population and sampling procedure}
The target population for the post-consolidation sample frame is the 20
commercial banks listed by Nepal Rastra Bank as of mid-October 2023
\citep{nrb2023list}. From this population, a
sample of 5 commercial banks is selected using a purposive (judgmental) sampling
technique. The selection criteria require: (i) representation of the prominent
commercial bank mergers that defined the early 2023 consolidation wave, (ii)
inclusion of an unmerged reference bank (NMB Bank) for descriptive context,
and (iii) continuous availability of audited annual financial statements,
NEPSE trading records, and accounting measures by their relevant observation
dates across the entire study period. Because the
sample is non-probabilistic and comprises 5 institutions, the empirical findings
are interpreted specifically for the sampled banks and serve as an indicative
benchmark rather than an automatic generalization across all 20 commercial
banks.

The sampled commercial banks, identified by their current legal corporate names,
are:
\begin{enumerate}
  \item Global IME Bank Limited (absorbed Bank of Kathmandu, joint operation commenced 9 January 2023, swap ratio 1:1) \citep{gbime2022agm,nrb2023list}
  \item Himalayan Bank Limited (acquired Civil Bank, joint operation commenced 24 February 2023, swap ratio 100:80.28) \citep{hbl2023report}
  \item Kumari Bank Limited (absorbed Nepal Credit and Commerce Bank, joint operation commenced 1 January 2023, swap ratio 1:1) \citep{kumari2022sgm,nrb2023list}
  \item Nepal Investment Mega Bank Limited (merger of Nepal Investment Bank and Mega Bank, joint operation commenced 11 January 2023, swap ratio 100:90) \citep{nimb2023report}
  \item NMB Bank Limited (maintained independent operations without merger during the period) \citep{nrb2023list}
\end{enumerate}
For each merged bank series, pre-merger observations represent the acquiring
institution and post-merger observations represent the consolidated
institution. The fiscal year in which a merger began joint operation is a
transition year: the audited statement's reporting boundary will be recorded
for each bank-year, and that observation will be flagged rather than described
as a clean pre-merger or post-merger year. The merger changes the scale and
composition of the series; regression coefficients are not interpreted as
causal estimates of merger effects. The study covers FY 2076/77--2081/82,
spanning years before, during, and after the early-2023 mergers, targeting 30
observations (5 bank series $\times$ 6 fiscal years), subject to the stated
data-availability criteria. The single unmerged series
provides descriptive context only; the model does not estimate a merger-versus-
non-merger comparison.

\msub{1.1.3 Nature and sources of data collection}
This study relies entirely on secondary financial and market data. Bank
financial statements provide the accounting measures; NEPSE trading records
provide market prices. NRB, SEBON, and Ministry of Finance publications
provide regulatory and sector context and are used for cross-checking where
they report the same observation. Sources for each bank-year are recorded in
the data file. The principal sources are:
\begin{itemize}
  \item Audited annual reports and published financial statements of the sampled commercial banks.
  \item \textit{Banking and Financial Statistics} and supervision reports published by Nepal Rastra Bank (NRB).
  \item Trading records, official share price bulletins, and annual publications of the Nepal Stock Exchange (NEPSE) and the Securities Board of Nepal (SEBON).
  \item \textit{Economic Survey} published by the Ministry of Finance (MoF), Government of Nepal.
\end{itemize}
Extracted figures are checked against the originating report and, where an
independent source reports the same measure, cross-checked before processing.
Differences are recorded and resolved using the originating audited report for
accounting measures and the NEPSE record for trading prices.

\msub{1.1.4 Research framework and definition of variables}
The research framework illustrating the functional relationship between the
explanatory accounting fundamentals and the dependent variable is depicted in
Figure~\ref{fig:framework}. Post-disclosure market price per share
($\mathit{MPS}^{post}$) serves as the dependent variable, while earnings per share
(EPS), dividend per share (DPS), book value per share (BVPS), and return on assets
(ROA) are the explanatory variables.

\begin{figure}[h]
\centering
\begin{tikzpicture}[node distance=0.45cm,
  ind/.style={rectangle, draw, minimum width=3.4cm, minimum height=0.7cm, text centered, font=\small},
  dep/.style={rectangle, draw, minimum width=3.6cm, minimum height=1.7cm, text centered, font=\small, fill=gray!10},
  arrow/.style={thick,->,>=stealth}]
  \node[ind] (eps)  {EPS};
  \node[ind, below=of eps] (dps)  {DPS};
  \node[ind, below=of dps] (bvps) {BVPS};
  \node[ind, below=of bvps] (roa) {ROA};
  \path (dps) -- (bvps) coordinate[midway] (mid);
  \node[dep, right=2.2cm of mid] (mps) {$\mathit{MPS}^{\mathit{post}}$};
  \draw[arrow] (eps.east)  -- ([yshift=18pt]mps.west);
  \draw[arrow] (dps.east)  -- ([yshift=6pt]mps.west);
  \draw[arrow] (bvps.east) -- ([yshift=-6pt]mps.west);
  \draw[arrow] (roa.east)  -- ([yshift=-18pt]mps.west);
  \node[above=0.2cm of eps, font=\small\itshape] {Independent variables};
  \node[below=0.25cm of mps, font=\small\itshape] {Dependent variable};
\end{tikzpicture}
\caption{Research framework (Source: adapted from Pradhan and Dahal, 2016)}
\label{fig:framework}
\end{figure}

Post-disclosure market price per share ($\mathit{MPS}^{\mathit{post}}$),
measured in Nepalese rupees (NPR), is the closing price on the first NEPSE
trading day after public disclosure of the corresponding audited annual
financial statements through NEPSE. The same rule is applied across all
sampled banks and fiscal years. The statement disclosure date, trading date,
and source records are logged for each bank-year. If a bonus issue, rights
issue, split, or other
share-capital action affects the quoted price between fiscal year-end and the
measurement date, the event and any price adjustment are documented using
NEPSE or issuer notices; adjusted and unadjusted prices are not mixed. Earnings
per share (EPS), in NPR per share, is profit attributable to ordinary equity
holders divided by the weighted-average ordinary shares outstanding, as
reported in or calculated from the audited statements. Dividend per share
(DPS), in NPR per share, is the amount reported for the corresponding fiscal
year; the data record identifies whether it is recommended, approved, or
distributed and records its public disclosure date. A DPS value that was not
publicly disclosed by the MPS observation date is recorded as unavailable and
is not used in the regression for that bank-year. Book value per share (BVPS), in
NPR per share, is the reported figure in the audited annual report. Return on
assets (ROA), in percent, is net profit after tax divided by year-end total
assets and multiplied by 100. For every bank-year, the data record identifies
the audited-report numerator and denominator and preserves any reported ROA
formula. ROA is recalculated from audited figures using the common definition
where both components are available. If a required figure cannot be
established, ROA is recorded as unavailable; a value using a different
denominator convention is not pooled with this measure.

\needspace{6\baselineskip}
\msub{1.1.5 Methods of analysis}
To examine the empirical relationships and evaluate the hypotheses formulated in
this study, both descriptive and inferential econometric techniques are applied.

\textbf{Model Specification:}
To examine the joint and conditional associations of all four accounting
fundamentals with post-disclosure market price per share, the following multiple
linear regression model is specified:
\begin{equation}
\mathit{MPS}^{post}_{it} = \beta_0 + \beta_1 \mathit{EPS}_{it} + \beta_2 \mathit{DPS}_{it} + \beta_3 \mathit{BVPS}_{it} + \beta_4 \mathit{ROA}_{it} + \varepsilon_{it}
\label{eq:main}
\end{equation}
where $\mathit{MPS}^{post}_{it}$ represents the closing market price observed
at the pre-specified post-disclosure trading date for bank $i$ in fiscal year
$t$. $\mathit{EPS}_{it}$, $\mathit{DPS}_{it}$, $\mathit{BVPS}_{it}$, and
$\mathit{ROA}_{it}$ denote earnings per share, dividend per share, book value
per share, and return on assets, respectively. The parameter $\beta_0$ represents
the intercept, $\beta_1$ through $\beta_4$ represent regression slope coefficients
capturing the estimated conditional association between each regressor and
$\mathit{MPS}^{post}$ holding the other included variables constant, and
$\varepsilon_{it}$ is the stochastic disturbance term.

\textbf{Statistical Tools and Estimation Procedure:}
The empirical analysis follows this protocol. Data suitability and model
assumptions are assessed before coefficient tests are interpreted:
\begin{enumerate}
  \item \textbf{Descriptive Statistics:} The mean, standard deviation, minimum,
    and maximum of all variables are computed to examine distributional
    properties across the panel.
  \item \textbf{Correlation Analysis:} Pearson correlations measure bivariate
    associations between each fundamental and $\mathit{MPS}^{post}$ and screen
    for collinearity among regressors.
  \item \textbf{Multiple Regression Analysis:} Eq.~\ref{eq:main} is
    estimated by pooled OLS as an exploratory baseline. The coefficient of
    determination ($R^2$) and adjusted $R^2$ describe in-sample fit. Individual
    coefficients are assessed with two-tailed $t$-tests at the 5\% level
    ($\alpha=0.05$), and joint model significance is assessed with an $F$-test
    at the 5\% level. Since the intended 30 observations repeat five bank series and
    four series cross merger boundaries, these conventional model-based tests
    are interpreted cautiously; they do not establish causal effects or
    sector-wide relationships.
  \item \textbf{Diagnostic Tests:} Model assumptions are examined before interpreting results.
    Multicollinearity is evaluated with variance inflation factors (VIF) and tolerance statistics alongside the correlation matrix.
    Residual diagnostics include checks for heteroskedasticity (White or
    Breusch-Pagan test), influential observations (Cook's distance), and
    residual distribution (Q-Q plot). The correlation matrix, VIF, and
    tolerance are reviewed for multicollinearity. Residual patterns are also
    inspected by bank series and fiscal year because repeated observations may
    be dependent. With only five bank series, bank-clustered standard errors
    would rely on very few clusters; conventional cluster-based tests can
    over-reject with small cluster counts \citep{cameron2008clusters}. They are
    therefore not treated as an automatic remedy. The small sample,
    transition-year reporting boundaries, and limits of conventional OLS
    inference are stated with the results. If diagnostics indicate serious
    violations, coefficient tests will not be presented as suitable for
    confirmatory interpretation.
\end{enumerate}

Data organization, coding, descriptive statistics, correlations, pooled OLS,
and the specified diagnostic checks are conducted using Microsoft Excel and
the Statistical Package for the Social Sciences (SPSS).
